\section{Related Work}

\subsection{Anterior Segment and Cataract Segmentation}

Encoder-decoder architectures have become the dominant paradigm for medical image segmentation since Ronneberger et al.\ introduced U-Net \cite{ronneberger2015unet}, demonstrating that a contracting path paired with a symmetric expanding path enables precise localization from very few annotated images. This architecture has been widely adopted in ophthalmic segmentation due to its ability to preserve fine spatial detail through skip connections.

The nnU-Net framework \cite{isensee2021nnunet} extended this foundation by automatically configuring the network architecture, preprocessing, and training strategy for any given dataset. Its strong performance on 23 public benchmarks established it as a competitive baseline for medical segmentation tasks, including ophthalmic structures, without task-specific engineering.

Surgical scene understanding in cataract procedures presents additional challenges due to specular reflections, occlusions from instruments, and intraoperative illumination changes. The CaDIS dataset \cite{grammatikopoulou2021cadis} addressed this gap by providing richly annotated frames of cataract surgery videos covering anatomical structures and surgical instruments. Building on such data, Ghamsarian et al.\ proposed DeepPyram \cite{ghamsarian2022deeppyram}, a pyramid-view architecture with deformable receptive fields designed to handle the multi-scale structure variability inherent in cataract surgery segmentation. DeepPyram demonstrated that standard convolutions benefit from explicit multi-scale context, motivating our use of complementary structural representations to further enrich feature diversity.

\subsection{Multimodal Fusion in Medical Imaging}

The transformer architecture introduced by Vaswani et al.\ \cite{vaswani2017attention} established attention mechanisms as a principled way to model long-range dependencies and cross-sequence interactions. Its key insight --- that queries from one representation can attend to keys and values from another --- directly motivates the cross-modal fusion strategy we adopt at the bottleneck.

In the medical imaging domain, Zhang et al.\ proposed mmFormer \cite{zhang2022mmformer}, a multimodal medical transformer that handles incomplete modality sets for brain tumor segmentation. By distributing modality-specific encoders connected through intra- and inter-modal transformers, mmFormer demonstrated that cross-modal attention outperforms simple feature concatenation when input representations are structurally heterogeneous. Similarly, Zhang et al.\ \cite{zhang2024dattnet} introduced DATTNet, a dual-attention transformer that models spatial and channel dependencies simultaneously in a hybrid CNN-transformer framework, achieving state-of-the-art results on multi-modal medical segmentation benchmarks.

Li et al.\ proposed DECTNet \cite{li2024dectnet}, a dual-encoder architecture combining a convolutional branch with a Swin Transformer branch, fusing their complementary representations via a feature fusion decoder. This direct structural precedent confirms the effectiveness of parallel encoders with heterogeneous inductive biases for medical segmentation. In contrast to DECTNet, which fuses different network types processing the same modality, our approach fuses two encoders processing structurally distinct representations (photometric and geometric) of the same image.

The multi-modal ophthalmic AI survey by Wang et al.\ \cite{wang2024multimodal} reviewed deep learning across fundus photography, OCT, OCTA, and clinical metadata, explicitly identifying cross-modal feature fusion as an open research direction with clinical impact. This motivates our exploration of intra-image multimodal fusion for cataract segmentation specifically.

\subsection{Foundation Models for Ophthalmic Segmentation}

Recent advances in foundation models have shifted the landscape of medical image segmentation. Ma et al.\ introduced MedSAM \cite{ma2024medsam}, fine-tuning the Segment Anything Model (SAM) on over 1.5 million medical image-mask pairs across 10 imaging modalities. While MedSAM achieves strong generalization, its scale --- requiring 20 A100 GPUs for training --- makes it impractical for deployment on small ophthalmic datasets without substantial computational infrastructure.

The challenge of adapting large models to specialized scenes was addressed by Chen et al.\ with SAM-Adapter \cite{chen2023samadapter}, which introduces lightweight task-specific adapters to guide SAM toward domain-specific features such as camouflage and shadow. While adapter-based methods reduce fine-tuning cost, they still inherit the full SAM encoder and are optimized for scenarios with large pre-training data available.

In ophthalmology specifically, Zhou et al.\ introduced RETFound \cite{zhou2023retfound}, a masked autoencoder pre-trained on 1.6 million retinal images that generalizes across multiple ophthalmic disease detection tasks. Similarly, Qiu et al.\ proposed VisionFM \cite{qiu2024visionfm}, a multi-modal multi-task vision foundation model validated on eight ophthalmic image modalities for disease diagnosis, prognosis, and anatomical segmentation.

\subsection{Positioning of the Proposed Work}

The above literature reveals three complementary trends: (1) task-specific architectures like U-Net and its variants achieve strong segmentation performance but are sensitive to the information available in a single imaging modality; (2) cross-modal attention mechanisms effectively fuse heterogeneous representations but have not been applied to intra-image geometric decomposition in cataract segmentation; and (3) foundation models provide generalization at scale but are inaccessible for small, specialized datasets.

Our proposed architecture bridges the gap between (1) and (2): we adapt the dual-encoder, cross-attention paradigm demonstrated in mmFormer and DECTNet to a lightweight, purpose-built model where the two encoders process the original RGB image and its Canny edge map respectively. This design explicitly provides the decoder with complementary photometric and geometric features through a bottleneck cross-attention module, without requiring large-scale pre-training. Compared to foundation model adaptation (3), this approach is computationally accessible and directly optimized for the structural characteristics of cataract anterior segment images.
